% =====================================================================
% Appendix: the two FunctionGraph-RAG selectors.
% Every number here is read from the shipped artifact files. Do not edit
% numbers by hand: pot_selector_frozen.json and cot_selector_frozen.json
% in the artifact carry them, and a reviewer can diff them against this.
% =====================================================================

\section{FunctionGraph-RAG: construction of both selectors}
\label{app:selector}

The two reasoning chains use separate selectors. They share a name and a purpose and almost
nothing else, so we describe them separately rather than presenting an average of the two. The
main-text results in Table~\ref{tab:rq2-full} are produced by the Program-of-Thought selector.

\begin{table}[h]
\centering
\small
\setlength{\tabcolsep}{4pt}
\begin{tabular}{@{}lll@{}}
\toprule
& \textbf{PoT chain} & \textbf{CoT chain} \\
& GPT-4o & DeepSeek-R1 \\
\midrule
Retrieval        & Contriever top-$30$   & BM25 top-$30$ \\
Graph edges      & shared title $\cup$ kNN & shared quantity \\
Expansion        & one hop, uncapped     & two hops, capped at $80$ \\
Selector         & linear, $4$ features  & linear, $56$ features \\
Injected         & at most $3$           & top $10$ \\
\addlinespace[2pt]
Labels           & $511$                 & $890$ \\
Splits used      & Easy, Medium          & Easy, Medium, Hard \\
Train/val/test   & $403$/$66$/$42$       & $678$/$106$/$106$ \\
Objective        & pairwise ranking      & pairwise ranking \\
\bottomrule
\end{tabular}
\caption{The two selectors. Both are fitted on direct \texttt{function\_id} relevance labels from
FinanceReasoning and neither sees a \textsc{FinExam-10K} item, answer or correctness signal at any
point.}
\label{tab:selectors}
\end{table}

\paragraph{Program-of-Thought chain.}
Candidates are the frozen Contriever top-$30$ cached from the upstream run, asserted at load time
to be exactly $30$ unique identifiers present in the function catalogue. Those $30$ are expanded by
one hop over a static graph whose edge set is the union of two relations: two functions are
adjacent when they share an \texttt{article\_title}, and each function is additionally linked to
its $k{=}4$ nearest neighbours under Contriever inner product. The graph is built once from the
function corpus alone, and the builder never opens a benchmark item, an answer, a model output or
a correctness flag. It has $3{,}133$ nodes and $15{,}889$ directed edges, degree minimum $4$,
median $5$, maximum $11$, and no isolated node. Expansion is exactly one hop and is not capped:
neighbours of the $30$ roots enter the pool, neighbours of neighbours do not. The resulting pool
has median size $98$.

A frozen linear model scores the pool over four features, the reciprocal of retrieval rank, the
reciprocal of graph depth, edge weight, and a degree scale, and the top three are injected. It was
fitted on $511$ direct \texttt{function\_id} relevance labels drawn from the Easy and Medium splits
of FinanceReasoning, $403$ for training, $66$ for validation, $42$ for test, under a pairwise
ranking objective, reaching test pair accuracy $0.860$. Of $752$ labelled examples, $241$ were
excluded because the labelled function does not appear in the candidate pool under this retrieval
protocol, which is recorded in the artifact rather than silently dropped.

\paragraph{Chain-of-Thought chain.}
Candidates are BM25 top-$30$. The graph is different in kind: two functions are adjacent when they
share a normalised quantity name, where a quantity is a function's Python parameter name or the
name of the variable it returns, matched after lowercasing, tokenisation, stopword removal, a
suffix singulariser and a nineteen-entry alias table. Expansion is nominally two hops with a cap of
$80$ candidates. In practice the cap is reached during the first hop on every query we inspected,
so the second hop contributes almost nothing, and the trained weights on the two-hop indicator
features are exactly zero. A linear reranker over $56$ features returns the top ten. It was fitted
on $890$ relevance labels from the Easy, Medium and Hard splits, $678$/$106$/$106$.

\paragraph{Which artifact produced the reported numbers.}
The upstream release contains a second Chain-of-Thought selector, fitted on Easy and Medium only,
with $736$ accepted labels and a $560$/$88$/$88$ split. It shares a \texttt{model\_version} string
with the one used here, because that field is a hardcoded constant rather than a real identifier,
and the two are distinguishable only by their file hashes. \textbf{The $736$-label selector did not
produce any number reported in this paper.} The results we report were run with the $890$-label
artifact, and both files ship so the distinction can be checked by hash.

\paragraph{Reproduction boundary.}
Selector inference is fully reproducible from the artifact: the graph, the coefficients, the
candidate protocol and the expansion code are all released. Selector \emph{training} is not, since
it consumes FinanceReasoning relevance labels that belong to that dataset. What is released instead
is the training code, the split counts, the accepted and rejected label counts, and a
candidate-protocol fingerprint that pins training-time and inference-time candidate construction to
the same definition, so a reader can verify that the selector was not fitted under one retrieval
protocol and deployed under another.

\paragraph{Relation to prior work.}
Our implementation follows the Function-RAG formulation rather than reproducing it exactly. The
selector is trained on FinanceReasoning and applied to \textsc{FinExam-10K} without adaptation,
which makes every number we report an out-of-distribution transfer. We state this rather than
claiming a faithful reimplementation, because the differences above are real and a reader
comparing against the original should know where they are.
